% area: Category Theory

\begin{definition}[Category] \label{def:category} 
A \emph{category} $\mathcal{C}$ consists of the following data:
\begin{enumerate} 
  \item A collection of objects, denoted $\operatorname{Ob}(\mathcal{C})$. 
  \item For every pair of objects $X,Y \in \operatorname{Ob}(\mathcal{C})$, a set $$ \operatorname{Hom}_{\mathcal{C}}(X,Y) $$ whose elements are called \emph{morphisms} (or arrows) from $X$ to $Y$. 
  \item For every triple of objects $X,Y,Z$, a composition map $$ \operatorname{Hom}_{\mathcal{C}}(Y,Z) \times \operatorname{Hom}_{\mathcal{C}}(X,Y) \to \operatorname{Hom}_{\mathcal{C}}(X,Z), $$ written $$ (g,f) \mapsto g \circ f . $$ 
  \item For every object $X$, an identity morphism $$ \operatorname{id}_X \in \operatorname{Hom}_{\mathcal{C}}(X,X). $$ 
  \end{enumerate} 
These data satisfy the following axioms: 
\begin{enumerate} 
\item \textbf{Associativity:} For morphisms $$ X \xrightarrow{f} Y \xrightarrow{g} Z \xrightarrow{h} W $$ we have $$ h \circ (g \circ f) = (h \circ g) \circ f . $$ 
\item \textbf{Identity:} For every morphism $f : X \to Y$, $$ f \circ \operatorname{id}_X = f \qquad\text{and}\qquad \operatorname{id}_Y \circ f = f . $$ \end{enumerate}
\end{definition}



\begin{remark}[Notation for Morphisms]
\label{rmk:morphism-notation}

Let $\mathcal{C}$ be a category as in Definition~\ref{def:category}.  
If $f \in \Hom_{\mathcal{C}}(X,Y)$ we write
$$
f : X \to Y.
$$

The object $X$ is called the \emph{domain} of $f$ and $Y$ the \emph{codomain}.

Composition of morphisms is written
$$
g \circ f : X \to Z
$$
whenever $f : X \to Y$ and $g : Y \to Z$ as in Definition~\ref{def:category}.
\end{remark}


\begin{definition}[Isomorphism]
\label{def:isomorphism}

Let $\mathcal{C}$ be a category (Definition~\ref{def:category}).  

A morphism
$$
f : X \to Y
$$
is called an \emph{isomorphism} if there exists a morphism
$$
g : Y \to X
$$
such that
$$
g \circ f = \operatorname{id}_X
\qquad
f \circ g = \operatorname{id}_Y.
$$

The morphism $g$ is called the \emph{inverse} of $f$ and is denoted
$$
f^{-1}.
$$

If such a morphism exists we say that $X$ and $Y$ are \emph{isomorphic} and write
$$
X \cong Y.
$$

This notion uses the identity morphisms and composition introduced in
Definition~\ref{def:category}.
\end{definition}


\begin{proposition}[Uniqueness of Inverses]
\label{prop:inverse-unique}

Let $\mathcal{C}$ be a category (Definition~\ref{def:category}).  
If a morphism $f : X \to Y$ is an isomorphism (Definition~\ref{def:isomorphism}),
then its inverse is unique.
\end{proposition}

\begin{proof}

Suppose
$$
g,h : Y \to X
$$
are both inverses of $f$. Then

$$
g \circ f = \operatorname{id}_X,
\qquad
f \circ g = \operatorname{id}_Y,
$$

and

$$
h \circ f = \operatorname{id}_X,
\qquad
f \circ h = \operatorname{id}_Y.
$$

Using associativity of composition from Definition~\ref{def:category} we obtain

$$
g
=
g \circ \operatorname{id}_Y
=
g \circ (f \circ h)
=
(g \circ f) \circ h
=
\operatorname{id}_X \circ h
=
h.
$$

Thus the inverse is unique.
\end{proof}


\begin{corollary}
\label{cor:iso-equivalence-relation}

Let $\mathcal{C}$ be a category (Definition~\ref{def:category}).  
The relation
$$
X \cong Y
$$
defined via isomorphisms (Definition~\ref{def:isomorphism})
is an equivalence relation on $\operatorname{Ob}(\mathcal{C})$.
\end{corollary}

\begin{proof}

\textbf{Reflexivity.}  
For every object $X$, the identity morphism $\operatorname{id}_X$ from
Definition~\ref{def:category} satisfies

$$
\operatorname{id}_X \circ \operatorname{id}_X = \operatorname{id}_X.
$$

Thus $\operatorname{id}_X$ is an isomorphism.

\textbf{Symmetry.}  
If $f : X \to Y$ is an isomorphism, Definition~\ref{def:isomorphism}
provides an inverse $f^{-1} : Y \to X$.

\textbf{Transitivity.}  
If $f : X \to Y$ and $g : Y \to Z$ are isomorphisms, then

$$
(g \circ f)^{-1} = f^{-1} \circ g^{-1}
$$

by associativity of composition in Definition~\ref{def:category}.
\end{proof}


\begin{definition}[Subcategory]
\label{def:subcategory}

Let $\mathcal{C}$ be a category (Definition~\ref{def:category}).  
A \emph{subcategory} $\mathcal{D}$ of $\mathcal{C}$ consists of

\begin{enumerate}

\item a collection of objects
$$
\operatorname{Ob}(\mathcal{D}) \subseteq \operatorname{Ob}(\mathcal{C})
$$

\item for each pair of objects $X,Y \in \operatorname{Ob}(\mathcal{D})$,
a subset

$$
\Hom_{\mathcal{D}}(X,Y)
\subseteq
\Hom_{\mathcal{C}}(X,Y)
$$

\end{enumerate}

such that

\begin{enumerate}

\item identity morphisms from Definition~\ref{def:category} belong to
$\mathcal{D}$,

\item composition in $\mathcal{D}$ is the restriction of composition in
$\mathcal{C}$.

\end{enumerate}

\end{definition}


\begin{definition}[Opposite Category]
\label{def:opposite-category}

Let $\mathcal{C}$ be a category (Definition~\ref{def:category}).  
The \emph{opposite category} $\mathcal{C}^{op}$ is defined as follows.

\begin{enumerate}

\item Objects:
$$
\operatorname{Ob}(\mathcal{C}^{op}) = \operatorname{Ob}(\mathcal{C})
$$

\item Morphisms:
$$
\Hom_{\mathcal{C}^{op}}(X,Y)
=
\Hom_{\mathcal{C}}(Y,X)
$$

\item Composition is defined by reversing composition in $\mathcal{C}$.

\end{enumerate}

Associativity and identity laws follow directly from those in
Definition~\ref{def:category}.
\end{definition}


\begin{remark}[Duality Principle]
\label{rmk:duality}

Let $\mathcal{C}$ be a category (Definition~\ref{def:category}).  
Any statement about $\mathcal{C}$ has a \emph{dual statement}
obtained by replacing

\begin{itemize}

\item morphisms $f : X \to Y$ by $f : Y \to X$

\item compositions $g \circ f$ by $f \circ g$

\end{itemize}

which corresponds precisely to passing to the opposite category
$\mathcal{C}^{op}$ defined in Definition~\ref{def:opposite-category}.

This observation is called the \emph{principle of duality}.
\end{remark}
